\documentclass[12pt]{article}
\usepackage[utf8]{inputenc}
\usepackage{amsmath, amssymb, amsthm}
\usepackage{geometry}

\geometry{a4paper, margin=1in}

\title{Challenge Problem Set 2}
\author{Math 292 Spring 2026 \\ Rutgers University}
\parskip=1ex
\begin{document}

\maketitle

(Basis) Since $8_c = 1(3_c) + 1(5_c)$,
S(8) is true.
Since $9_c = 3(3_c)$
S(9) is true.
Since $10_c = 2(5_c)$
S(10) is true.

(Inductive Step) Let $k \geq{10}$ be an arbitrary natural number.
Assume S(8), S(9), S(10), ... , S(k) are all true.

By inductive hypothesis, there are non-negative integers $a$ and $b$ such that
\begin{align*}
    (k-2)_c = a(3_c) + b(5_c)
\end{align*}

Example: $183k +19l = gcd(183, 19)$

Euclidean Algorithm:
\begin{align*}
    183 & = 19*9 + 12  \\
    19  & = 12 * 1 + 7 \\
    12  & = 7 * 1 + 5  \\
    7   & = 5 * 1 + 2  \\
    5   & = 2 * 2 + 1  \\
    1   & = 1 * 1 + 0
\end{align*}

Recursive:
\begin{align*}
    1 & = 1(5) - (2)2         \\
      & = 1(5) - 2(7-5)       \\
      & = 3(5) - 2(7)         \\
      & = 3(12-7) - 2(7)      \\
      & = 3(12) - 5(7)        \\
      & = 3(12) -5(19-12)     \\
      & = 8(12) -5(19)        \\
      & = 8(183-19(9)) -5(19) \\
      & = 8(183) -77(19)
\end{align*}

Thus, $k = 8$ and $l = -77$.

\newpage{}

Suppose, for the sake of contradiction, there are no such subsequences.
That is, there is no increasing subsequences of length 4 or more,
and there is no decreasing subsequences of length 4 or more.

For each $i \in{A}$, let $I_i=$ largest length of an increasing subsequences starting with $i$, and
$D_i$ = largest length of a decreasing subsequences starting with $i$.

So, $I_i, D_i \in{\mathbb{N}}$.

Then, $1 \leq{} I_i, D_i \leq 3$ for each $i \in A$.

So, we have $3(3)$ possible distinct $(I_i, D_i)$ pairs.

Let these 9 distinct pairs be the pigeonholes and the 10 numbers be the pigeons.

So, by PHP, some two distinct numbers, say, $i$ and $j$, will be assigned to the same pair.

So, $(I_i, D_i) = (I_j, D_j)$

Assume, without loss of generality, $i < j$.

Case 1. Suppose $i$ appears before $j$ in the rearrangement.

Then, there is an increasing subsequence starting at $j$ w/ length $I_j$.

But, $I_i = I_j$. So that, subsequence along i creates an subsequence of length $I_i + 1$.

This is contradicts that the $I_i$ is the longest subsequence.

\vspace{15pt}
Case 2. Suppose $j$ appears before $i$ in the rearrangement.

Then, there is an decreasing subsequence starting at $i$ w/ length $D_i$.

But, $D_i = D_j$. So, that subsequence along j creates an subsequence of length $D_j + 1$.

This is contradicts that the $D_j$ is the longest subsequence.

\newpage{}

Claim: 32 is the largest number of red squares that works.

Divide the board into 2x2 squares with sides parallel to the board.

Let the 16 2x2 squares be the pigeonholes and the 33 red 1x1 squares be the pigeons.
\end{document}

